\chapter{典型计算内核模式}

先看一个很容易被忽略的事故。团队要统计日志中最活跃的一百个用户，于是把输入切成多个 chunk，每个 worker 在自己的 chunk 中维护局部 top-k，最后只合并这些局部候选。小数据测试通过，线上却漏掉了一个总量很高的用户。这个用户在每个 chunk 中都排在第一百名之外，所以局部 top-k 阶段把它丢掉了；全局合并阶段再也没有机会看见它。

这个错误不是某个堆操作写错了，而是内核组合的合同错了。局部 top-k 丢弃了后续精确全局 top-k 仍然需要的信息。单个内核在局部语义下正确，不代表它放进流水线后仍然正确。本章要训练的正是这种判断力：看到一个计算任务时，先问输入和输出的关系、输出位置如何确定、是否有共享写、数据访问是否规律、是否允许丢弃信息，然后再选择 map、reduce、scan、filter、partition、histogram、join、sort、top-k 或 graph frontier。

这里的 \term{内核}{kernel} 指一段高频、边界清楚、可以独立测量的计算片段。它不一定是操作系统内核，也不一定是 GPU kernel；在本册中，解析一批日志字段、过滤有效记录、统计事件类型、合并 partial、生成 top-k 候选，都可以是计算内核。每个内核必须有 \term{reference}{朴素但可信的对照实现}，用来固定语义；优化版本先和 reference 对齐，再谈吞吐。

\section{先写内核合同，而不是先写并行代码}

典型计算内核不能写成算法名清单。一个名字背后可能有完全不同的系统问题。Histogram 可以是 256 个小桶，也可以是亿级用户 ID 的稀疏聚合；filter 可以只输出计数，也可以要求稳定保留原始顺序；join 可以是一张小表查大表，也可以是两张大表产生巨大输出。若不写清输入形状和输出语义，任何“最佳实现”都没有意义。

\term{合同}{contract} 是内核对上下游作出的承诺。它至少回答九个问题：输入字段是什么，输出字段是什么，是否保持顺序，输出位置是否天然确定，是否丢弃信息，是否精确，是否可重试，是否有外部副作用，是否需要稳定的 tie-break 或随机种子。合同越早写清，后面的并发、I/O 和分布式章节越容易复用这一段设计。

\begin{lstlisting}[language=C++,caption={KernelContract 的教学形态}]
enum class KernelOutputShape {
  OneToOne,
  ZeroOrOne,
  ManyToOne,
  OneToMany,
  Reordered,
  Approximate
};

struct KernelContract {
  KernelOutputShape output_shape = KernelOutputShape::OneToOne;
  bool preserves_input_order = true;
  bool writes_shared_state = false;
  bool may_drop_information = false;
  bool requires_stable_tiebreak = false;
  bool materializes_output = false;
};
\end{lstlisting}

这段结构不是要求读者立刻实现框架，而是把思考顺序固定下来。Map 多数是一对一输出；filter 是一对零或一；reduce 是多对一；join 可能一对多；sort 和 partition 会重排；sketch 和近似 heavy hitter 可能丢弃信息。合同让这些差别变成代码边界，而不是藏在函数名里。

\begin{keyidea}
分析计算内核时，先问四件事：每个输入产生多少输出，输出位置是否已知，多个输入是否会写同一状态，内核是否会丢弃后续仍需要的信息。
\end{keyidea}

\topic{材料化、流水线和 pipeline breaker}
\term{材料化}{materialization} 是把中间结果写成数组、文件或其他可复查对象。它增加内存或 I/O 成本，却换来调试、复用、恢复和阶段测量。\term{流水线}{pipeline} 则让上游输出直接进入下游，减少中间写入，但背压、生命周期和错误传播更复杂。

\term{pipeline breaker}{会打断流水线的阶段} 是必须等待一批数据、重排数据、聚合状态或持久化边界后才能继续的内核。Sort、global reduce、join build、shuffle、checkpoint 常常是 breaker；简单 map 和普通 filter 更容易流水线。识别 breaker 不是运行时的附加工作，而是内核合同的一部分。

\topic{manifest 让中间结果可复查}
当 partition 或 shuffle 生成多个输出块时，只返回一个 \code{std::vector} 不够。系统还需要知道每个块属于哪个 bucket、包含多少记录、写到哪里、校验和是什么、是否已经提交。\term{manifest}{清单} 就是描述这些输出对象身份和边界的记录。单机 partition 的 manifest 到第 18 章会自然变成分布式 shuffle block 的 manifest。

\section{一对一和多对一：Map 与 Reduce}

Map 是最容易理解的内核：第 \code{i} 个输入通常产生第 \code{i} 个输出。它没有跨元素依赖，适合顺序扫描、SIMD 和线程切分。可是 map 也最容易被低估。若每个元素只做一次乘法和一次加法，它很快会受内存带宽限制；若每个元素调用复杂函数、访问冷字段或分配字符串，瓶颈又会转移到前端、分支或分配器。

\begin{lstlisting}[language=C++,caption={一对一 map 的 reference}]
void scale_and_shift(std::span<const float> input,
                     std::span<float> output,
                     float scale,
                     float shift) {
  const std::size_t count = std::min(input.size(), output.size());
  for (std::size_t index = 0; index < count; ++index) {
    output[index] = input[index] * scale + shift;
  }
}
\end{lstlisting}

这个 reference 的正确性来自输出位置合同：每个输出位置只依赖同一位置输入。优化时可以切块、向量化或融合相邻 map，但这些优化不能改变位置关系。若输入输出区间可能破坏性重叠，还要把别名语义写进接口；否则编译器和维护者都无法确定循环能否安全重排。

\topic{Map 融合减少字节，不是为了让代码更复杂}
连续三个 map 若每一步都写出中间数组，会产生多轮内存读写。融合可以把这些逐元素变换放进一次循环，减少中间材料化。可是融合后循环体更大，寄存器压力更高，错误路径也可能混进热循环。融合是否值得，要比较未融合耗时、融合耗时、中间数据大小、错误报告是否变难，而不是凭直觉判断。

日志项目中，timestamp 归一化、event type 编码、用户 ID 映射都像 map。若解析器输出对象数组，map 会拖着冷字段一起走；若解析器输出列式 \code{HotBatch}，map 可以只扫热字段。第 6 章的数据布局因此不是孤立优化，而是 map 内核能否真正顺序扫描的前提。

\topic{Reduce 的第一问题是合并语义}
Reduce 把多个输入合成一个输出。它看起来像“加起来”，但真正的问题是合并操作的语义。整数最大值有清楚的单位元和结合律；浮点求和不严格结合，合并树会影响舍入；字符串拼接可能满足语义结合，却被分配成本支配；错误列表合并可能要求稳定顺序。没有合并语义，就没有正确的并行 reduce。

并行 reduce 的反例，是所有线程直接更新一个共享结果。下面的代码可能得到正确计数，但每个命中元素都要修改同一条 cache line，扩展性会很差。

\begin{lstlisting}[language=C++,caption={反例：高频共享写的 reduce}]
void count_positive_bad(std::span<const std::int32_t> values,
                        std::atomic<std::int64_t>& total) {
  for (const std::int32_t value : values) {
    if (value > 0) {
      total.fetch_add(1, std::memory_order_relaxed);
    }
  }
}
\end{lstlisting}

\code{memory_order_relaxed} 只放松了顺序约束，并没有让共享写消失。更可靠的起点是每个 worker 写自己的 partial，最后合并。

\begin{lstlisting}[language=C++,caption={线程本地 partial，把高频共享写移出热路径}]
constexpr std::size_t CACHE_LINE_BYTES = 64;

struct alignas(CACHE_LINE_BYTES) CountPartial {
  std::int64_t positive = 0;
};

void count_positive_range(std::span<const std::int32_t> values,
                          CountPartial& partial) {
  std::int64_t local_count = 0;
  for (const std::int32_t value : values) {
    if (value > 0) {
      ++local_count;
    }
  }
  partial.positive = local_count;
}
\end{lstlisting}

这个版本的原则比代码更重要：热路径写寄存器和本地 partial，不写全局共享对象。它把问题从“每个元素一次同步”变成“每个 worker 一次合并”。当 partial 很大时，合并本身也要继续设计，可能需要树形合并、按 NUMA 节点合并或按 bucket 并行合并。

\section{输出位置未知：Scan、Filter 与 Partition}

Filter 选择满足谓词的元素。串行版本可以直接 \code{push_back}，因为当前输出位置由前面已经保留了多少元素决定。并行后，这个隐含状态变成共享资源。多个线程同时 \code{push_back} 会争用；使用原子下标可以分配位置，却可能破坏稳定顺序；加锁保持顺序又回到瓶颈。

稳定 filter 的推导路线是三阶段。第一，每个块统计自己会保留多少元素。第二，对块计数做 scan，得到每个块的全局输出起点。第三，每个块再次扫描输入，写入自己的连续输出区间。Scan 在这里不是抽象算法题，而是输出位置分配器。

\begin{lstlisting}[language=C++,caption={过滤前先统计每个块的输出数量}]
std::int64_t count_selected(std::span<const std::int32_t> values,
                            std::int32_t threshold) {
  std::int64_t selected_count = 0;
  for (const std::int32_t value : values) {
    if (value >= threshold) {
      ++selected_count;
    }
  }
  return selected_count;
}
\end{lstlisting}

这个函数只统计，不写输出，因此没有共享写。对所有块的 \code{selected_count} 做前缀和，就能得到每块的输出起点。后续写出时，每个 worker 都只写自己负责的区间。它多读了一遍输入，却消除了共享 push、动态扩容和顺序混乱。

\topic{选择率决定 filter 的成本模型}
Filter 的性能高度依赖选择率。选择率为 0 时，输出很小，但仍要读完整输入；选择率为 100\% 时，filter 接近 copy；选择率在中间且随机时，分支和输出位置都复杂。低选择率可能适合 selection vector；高选择率可能适合 bitmap 或直接保留原数组加 mask；记录很大时，先保存 record id 而不是复制整条记录可能更好。

测试 filter 不能只测 50\% 随机选择。至少要覆盖全不通过、全通过、低选择率、高选择率、长记录、错误记录和最后一块不满。否则实现很容易只对某个友好输入成立。

\topic{Partition 是多路输出位置分配}
Partition 把元素按 bucket 或 key 放到不同区域。二路 partition 是 filter 的近亲，多路 partition 则是 shuffle 的前置形态。可靠做法仍然是先统计每个 worker 到每个 bucket 的数量，再做二维前缀和，最后让每个 worker 写入每个 bucket 中自己的区间。

分区不只是为了并行写输出，也可以为了改变后续访问形态。把记录按 event type 分桶，后续每个桶里的分支更稳定；把记录按 user id 分区，后续 group-by 的工作集更小；把记录按 reduce target 写成 block，分布式阶段就可以直接传输。分区的成本是额外扫描、输出写入和可能打乱顺序；收益是减少共享写、控制工作集或形成 stage 边界。

\section{共享写和数据倾斜：Histogram 与 Top-k}

Histogram 是共享写问题的典型训练场。它的计算很短：读 key，找到 bucket，计数加一。系统成本却取决于 bucket 数量、key 分布、线程数和内存布局。桶很少时，全局原子会形成热点；桶很多时，每线程完整桶数组可能占用大量内存；key 极度倾斜时，一个热点 bucket 会支配合并阶段。

\begin{lstlisting}[language=C++,caption={稠密 key 的本地直方图}]
constexpr std::int32_t EVENT_BUCKET_COUNT = 1024;

struct LocalHistogram {
  std::array<std::int64_t, EVENT_BUCKET_COUNT> buckets{};
};

void histogram_dense(std::span<const std::int32_t> event_types,
                     LocalHistogram& histogram) {
  for (const std::int32_t event_type : event_types) {
    if (0 <= event_type && event_type < EVENT_BUCKET_COUNT) {
      ++histogram.buckets[static_cast<std::size_t>(event_type)];
    }
  }
}
\end{lstlisting}

这个版本速度可能很高，因为 bucket 连续、无哈希、无共享写。但它要求 key 是小范围整数。若 key 是 64 位用户 ID，直接分配巨大数组不现实；若 key 是字符串，哈希、比较、分配和字典编码都会进入成本。工程中常见的路线是：小范围 key 用数组桶，稀疏 key 用本地哈希表或排序聚合，热点 key 用采样、隔离或二级聚合。

\topic{哈希聚合和排序聚合没有固定胜负}
Group-by 聚合通常有两条路线。哈希聚合边读边更新表，平均复杂度好，适合无序输入和点更新；排序聚合先按 key 排序，再顺序扫描相同 key，适合需要有序输出、key 可比较、内存局部性重要或哈希冲突严重的场景。排序聚合多了排序成本，却可能减少随机访问和分配。

选择时要看 key 数量、输入大小、输出是否需要排序、内存预算、倾斜程度和是否已有顺序。若输入已经按 key 近似有序，排序聚合可能很有吸引力；若 key 几乎全唯一，局部 combine 的收益很小；若 key 先做字典编码能变成稠密整数，数组桶可能比通用哈希表更合适。

\topic{Top-k 不能随便提前丢信息}
Top-k 的目标是减少数据移动：只保留最有可能进入最终结果的候选，而不是全量排序。每个 worker 维护局部小堆，最后合并候选，是常见结构。但它只有在合同允许时才正确。若局部阶段按 chunk 丢弃了全局精确结果仍然需要的 key，就会重演本章开头的事故。

精确全局 top-k 通常要先完成按 key 的完整聚合，再从全局计数中选择 top-k。若 key 基数太大，可以使用近似 heavy hitter 或更大的局部候选集，但必须写清误差合同：可能漏报什么，是否做二次精确验证，输出是否标记为近似。监控趋势可以接受近似，计费和审计通常不能。

\begin{lstlisting}[language=C++,caption={确定性 top-k 候选的接口形态}]
struct TopCandidate {
  std::int32_t key = 0;
  std::int64_t score = 0;
};

bool comes_before(const TopCandidate& left,
                  const TopCandidate& right) {
  if (left.score != right.score) {
    return left.score > right.score;
  }
  return left.key < right.key;
}
\end{lstlisting}

Tie-break 必须成为接口语义。若两个 key 分数相同，输出按 key 升序还是任意？不确定顺序会让测试、缓存和分布式重试难以复现。一个高质量 top-k 内核既要快，也要把稳定性写进合同。

\section{数据形状改变策略：Sort、Join、Stencil 与 Graph Frontier}

Sort 不只是算法题，也是系统工具。排序可以把相同 key 聚到连续区间，便于聚合、join、压缩和确定性输出。外部排序先生成有序 run，再多路归并；分布式排序先采样选择分区边界，再 shuffle 到 reducer。排序常常是 pipeline breaker，但它换来顺序访问、稳定输出和更容易复查的阶段边界。

\topic{Join 的成本在输出规模和构建侧}
Join 最容易低估的是输出放大。两个输入各一百万行，若某个 key 在左表出现一千次，在右表也出现一千次，该 key 的 join 输出就是一百万行。即使输入不大，输出也可能直接把内存打爆。Join 内核必须估计 key 分布、匹配度、输出上界和背压策略。

Hash join 通常把较小表构建成哈希表，再扫描较大表 probe；sort-merge join 先按 key 排序，再线性合并；nested loop 只适合极小输入或有强索引过滤的场景。选择不是固定的。若小表 key 极端倾斜，哈希 bucket 会很长；若大表过滤后很小，先 filter 再 join 更好；若最终输出需要按 key 有序，sort-merge join 可能减少后续成本。

\topic{Stencil 从邻域复用开始}
Stencil 每个输出读取邻域输入，例如一维平滑、二维卷积和网格有限差分。它的关键是复用邻域数据，避免重复从内存加载。内部区域和边界区域最好分开：内部热循环没有复杂分支，边界单独处理；或者使用 padding 和 ghost cell 统一访问规则。

并行 stencil 还要处理块之间的 halo 数据。单机上，halo 是相邻块共享的边界输入；分布式上，它会变成节点之间交换的消息。这个例子说明，cache blocking 和分布式 halo exchange 是同一个思想在不同尺度上的表现。

\topic{Graph frontier 暴露不规则访问}
图遍历是规则数组内核的反面。邻接表访问可能随机，顶点度数差异巨大，frontier 大小每层变化，写 visited 状态可能有竞争。CSR 把所有边放进连续数组，能改善邻接扫描；但顶点属性访问、frontier 去重和负载均衡仍然复杂。

BFS 的 top-down 策略从当前 frontier 扩展邻居；bottom-up 策略从未访问顶点检查是否有邻居在 frontier。当前沿很小时 top-down 更好，当前沿很大时 bottom-up 可能减少边访问。方向优化的本质，是根据数据形状切换内核，而不是固定套一个并行模板。

\begin{mentalmodel}
图遍历把本章的问题集中在一起：随机读、共享写、负载倾斜、输出去重、frontier 表示切换和任务调度。读懂图内核，后面看任务运行时的 ready set 会更自然。
\end{mentalmodel}

\section{把内核组合成可测流水线}

真实系统很少只运行一个内核。日志统计可以拆成一条流水线：parse map 把原始文本变成字段列；valid filter 生成选择向量和错误列表；event histogram 对有效记录做本地计数；partition 把 partial 或记录按 reduce target 分桶；merge reduce 合并所有 worker 的 partial；top-k 输出热门事件。每一步都有独立 reference，也有组合后的不变量。

组合后的验收比单内核更严格。Filter 后记录数加被过滤数应等于输入数；histogram 所有桶计数总和应等于 filter 输出数；partition 每个 key 的目标必须和 reduce 路由一致；top-k 的候选策略不能丢失精确结果需要的 key；最终 checksum 要和串行 reference 对齐。若单步都正确但组合失败，通常是顺序、身份、材料化边界或信息丢弃出了问题。

\topic{沿一条记录追踪}
调试组合流水线时，选一条记录从输入追到输出。它是否通过 filter，进入哪个 partition，更新哪个局部桶，在 join 中匹配几条维表记录，最后是否进入 top-k 候选。沿单条记录追踪，比只看每个函数的单元测试更容易发现接口错误，例如 filter 后顺序改变导致 join 假设失效，或者局部 combine 丢失错误报告所需的 record id。

\topic{规模账本比单点耗时更重要}
每个 stage 都要记录 records in、records out、bytes out、temporary bytes、state bytes、最大 bucket、join 放大倍数、错误数和 checksum。局部内核变快不代表系统变快。若 filter 选择率突然升高，join 输出会膨胀；若 partition 出现热点桶，下一阶段会长尾；若 top-k 为了保险保留太多候选，内存带宽可能被中间数据吃掉。

\begin{lstlisting}[language=C++,caption={阶段报告的最小字段}]
struct StageReport {
  std::string_view stage_name = {};
  std::uint64_t records_in = 0;
  std::uint64_t records_out = 0;
  std::uint64_t temporary_bytes = 0;
  std::uint64_t state_bytes = 0;
  std::uint64_t checksum = 0;
};
\end{lstlisting}

这不是完整观测系统，却表达了第二册的实验纪律：每次优化都要能回答输入规模、输出规模、临时数据、状态大小和结果校验。没有这些字段，吞吐数字很容易误导读者。

\section{项目落地：KernelPatternSuite}

Compute Systems Engine 在本章应形成一个小而完整的 \code{KernelPatternSuite}。它不需要一次写成通用框架，但必须覆盖几种不同形状：一对一 map、输出未知的 stable filter、共享写明显的 histogram、局部候选的 top-k、会输出放大的 join 或替代案例、frontier 或不规则负载的小实验。每个内核都保留 reference、优化版本、输入生成器、校验器和 stage report。

\topic{输入族要攻击假设}
压力输入不是简单放大规模，而是攻击内核假设。Filter 要测试全通过、全不通过、低选择率和高选择率；histogram 要测试均匀 key、Zipf 分布和单热点 key；top-k 要测试“分散强 key”，即每个 chunk 中不进局部 top-k、全局却很强的 key；join 要测试重复 key 导致的输出放大；graph 要测试规则网格图、随机图和幂律图。

\topic{源码阅读入口}
Linux 源码阅读可以从 \filepath{lib/sort.c}、\filepath{include/linux/rhashtable.h} 和若干 bitmap/hashtable 辅助实现附近进入。阅读目标不是把整套内核代码抄进书里，而是追一条和本章实验相关的路径：排序怎样处理元素移动，哈希表怎样处理扩容和桶，bitmap 怎样表达集合。读源码报告要区分证据等级：直接测得的数字、从控制流推导的解释、当前设备无法验证的合理猜测。

\begin{labbox}
实验：扩展事件流高性能管线。第一步实现线程本地 histogram，对比全局 mutex、全局 relaxed atomic 和本地 partial。第二步实现 stable filter，用块计数和 scan 分配输出位置。第三步加入 top-k，构造分散强 key，证明局部 top-k 不能作为精确全局 top-k 的前置。报告必须写 records in、records out、temporary bytes、state bytes、最大 bucket、checksum 和未验证边界。
\end{labbox}

\section{复查清单}

\topic{一、reference 是否仍然存在}
没有 reference，后面的 SIMD、多线程、异步或分布式改造都无法证明语义没有变。Reference 可以慢，但必须简单、边界清楚、错误路径可解释。

\topic{二、内核合同是否写清}
每个内核都要写输入、输出、顺序、精确性、信息丢弃、外部副作用、可重试性和 tie-break。若合同写不出来，先不要优化。

\topic{三、压力输入是否攻击核心边界}
只扩大数据量不够。要构造低选择率、高选择率、热点 key、稀疏 key、分散强 key、输出放大、frontier 长尾和边界长度输入。

\topic{四、状态是否可观察}
报告中应能看到关键对象的身份、大小、状态、错误数和校验结果。对 partition 而言，要有 bucket 范围；对 shuffle 而言，要有 manifest；对 top-k 而言，要有候选数量和 tie-break。

\topic{五、组合不变量是否成立}
单个内核正确不代表流水线正确。Filter 输出数、histogram 总和、partition 路由、join 输出规模、top-k 候选覆盖和最终 checksum 都要对齐 reference。

\topic{六、是否说明不采用更复杂方案的理由}
没有证据证明队列是瓶颈，就不应把锁队列换成无锁；没有证明输出顺序可变，就不应随意重排；没有证明错误路径可恢复，就不应只提交正常路径。能解释不做什么，说明工程判断开始成形。

\section{本章习题}

\begin{exercisebox}
习题一：把一个日志统计任务拆成 parse map、valid filter、partition、histogram、top-k 五个阶段。每个阶段写出输入、输出、是否保持顺序、是否有共享写、是否可能丢弃信息和主要瓶颈。

习题二：设计 stable filter。要求先写串行 \code{push_back} reference，再写共享锁反例，最后写块计数、scan、写出三阶段方案。测试覆盖选择率为 0、选择率为 100\%、随机选择、最后一块不满和错误记录。

习题三：为 histogram 选择三种策略：全局原子、本地桶、稀疏本地 map。分别说明适用的 bucket 数量、key 分布、线程数、内存占用和合并成本。

习题四：构造本章开头的 top-k 反例。让一个 key 在每个 chunk 中都不进入局部 top-k，但全局总和进入前 k。说明怎样修改流水线才能得到精确结果，怎样声明近似结果才算诚实。

习题五：为一个 join 任务估算输出放大。给定每个 key 在左右输入中的出现次数，计算最坏输出规模，并说明何时应选择 hash join、sort-merge join、预过滤或热点 key 单独处理。

习题六：用一个小图实验比较 vector frontier 和 bitmap frontier。记录每层 frontier 大小、边访问数、重复发现次数和 worker 空闲时间，说明何时切换表示。
\end{exercisebox}
